
\documentclass[12pt,a4paper]{article}

\usepackage[a4paper,margin=1in]{geometry}
\usepackage{graphicx}
\usepackage{booktabs}
\usepackage{amsmath,amssymb}
\usepackage{float}
\usepackage{hyperref}
\usepackage{xcolor}
\usepackage{listings}
\usepackage{enumitem}
\usepackage{fancyhdr}
\usepackage{longtable}

\hypersetup{colorlinks=true,linkcolor=blue,urlcolor=blue}

\pagestyle{fancy}
\fancyhf{}
\lhead{ICS1512}
\rhead{Machine Learning Algorithms Laboratory}
\cfoot{\thepage}

\lstset{
language=Python,
basicstyle=\ttfamily\small,
numbers=left,
frame=single,
breaklines=true,
showstringspaces=false
}

\begin{document}

\begin{tabular}{|p{4cm}|p{11cm}|}
\hline
Experiment No. & 6 \\ \hline
Experiment Title & Bagging, Boosting, and Stacked Ensemble Models\footnotemark \\ \hline
Student Name & Sanjay S \\ \hline
Register Number & 3122247001056 \\ \hline
Faculty & Dr. Poreddy Ajay Kumar Reddy \\ \hline
Submission Date & 21/09/2026 \\ \hline
\end{tabular}
\footnotetext{\textbf{GitHub Repository: }\url{https://github.com/sanjay-6727/Introdution-to-Machine-Learning-}}

\section{Objective}
Get a real feel for the three big ensemble strategies, Bagging, Boosting and Stacking, by actually
building one of each rather than just reading about them. Specifically: implement a Bagging classifier
and a couple of Boosting classifiers, put together a Stacked ensemble out of a few different base
learners, and then compare all three on accuracy, stability and how well they generalize, with an eye
on what each one is doing to bias and variance under the hood.

\section{Problem Statement}
Given the Breast Cancer Wisconsin Diagnostic dataset, preprocess it, split it 80/20, and build three
ensemble models: a Bagging classifier over Decision Trees, a Boosting classifier (both AdaBoost and
Gradient Boosting were tried), and a Stacked ensemble combining SVM, Naive Bayes and Decision Tree with
a Logistic Regression meta-learner. Every model gets its hyperparameters tuned with 5-fold
cross-validation, and the final comparison looks at accuracy, precision, recall and F1 on the same
held-out test set.

\section{Dataset Description}
\begin{longtable}{|p{4.2cm}|p{4.5cm}|p{1.6cm}|p{1.8cm}|p{2cm}|}
\hline
\textbf{Dataset} & \textbf{Source} & \textbf{Samples} & \textbf{Features} & \textbf{Target Classes} \\ \hline
Breast Cancer Wisconsin (Diagnostic) & UCI Machine Learning Repository, loaded via
\texttt{sklearn.datasets.load\_breast\_cancer} & 569 & 30 numeric & 2 (Malignant, Benign) \\ \hline
\end{longtable}

\textbf{Class distribution:} 357 Benign, 212 Malignant, so it's a bit skewed toward the benign class
but nowhere near severe enough to need SMOTE or class weighting, especially once I stratified the
train/test split on the label.

\textbf{Preprocessing steps:}
\begin{itemize}
\item Loaded the built-in scikit-learn version of the dataset, which is the same 569-sample, 30-feature
WDBC data as the UCI copy, just already cleaned and numeric.
\item No missing values to deal with here, so no imputation was needed.
\item An 80/20 train-test split, stratified on the class label (455 train, 114 test).
\item All 30 features standardized with \texttt{StandardScaler}, fit only on the training fold inside
each pipeline so it gets refit correctly during cross-validation.
\end{itemize}

\section{Ensemble Models Used}

\textbf{Bagging:} Decision Tree base estimator, tuned over the number of estimators, the fraction of
samples drawn per estimator, and the fraction of features drawn per split.

\textbf{Boosting:} Tried both AdaBoost and Gradient Boosting, each tuned over number of estimators and
learning rate (plus max depth for Gradient Boosting), to see which sequential-correction approach
handled this dataset better.

\textbf{Stacked Ensemble:} Base learners were an SVM (RBF kernel), Gaussian Naive Bayes and a Decision
Tree, feeding into a Logistic Regression meta-learner trained with 5-fold internal blending so the
meta-learner never sees predictions from data its own base learner was trained on.

\section{Hyperparameter Evaluation Results}

\begin{table}[H]
\centering
\caption{Bagging -- Hyperparameter Evaluation (top configurations by CV accuracy)}
\begin{tabular}{ccc}
\toprule
n\_estimators & max\_samples & Avg CV Accuracy (\%) \\
\midrule
50 & 0.5 & 96.04 \\
50 & 1.0 & 96.04 \\
100 & 0.7 & 96.04 \\
100 & 1.0 & 96.04 \\
50 & 0.7 & 95.82 \\
100 & 0.7 & 95.82 \\
\bottomrule
\end{tabular}
\end{table}

Best Bagging setting: \texttt{n\_estimators=50}, \texttt{max\_samples=0.5}, \texttt{max\_features=0.5}.
5-fold CV F1 on the training set: [0.965, 0.964, 0.946, 0.974, 0.991], mean = 0.9681.

\begin{table}[H]
\centering
\caption{Boosting -- Hyperparameter Evaluation (top AdaBoost configurations by CV accuracy)}
\begin{tabular}{ccc}
\toprule
n\_estimators & learning\_rate & Avg CV Accuracy (\%) \\
\midrule
100 & 1.0 & 98.02 \\
200 & 1.0 & 97.80 \\
200 & 0.1 & 97.14 \\
50 & 1.0 & 96.70 \\
100 & 0.1 & 96.04 \\
50 & 0.1 & 95.38 \\
\bottomrule
\end{tabular}
\end{table}

Best AdaBoost setting: \texttt{n\_estimators=100}, \texttt{learning\_rate=1.0}.
5-fold CV F1: [0.991, 1.000, 0.957, 0.983, 0.991], mean = 0.9844.

\begin{table}[H]
\centering
\caption{Gradient Boosting -- Hyperparameter Evaluation (top configurations by CV accuracy)}
\begin{tabular}{ccc}
\toprule
n\_estimators & learning\_rate & Avg CV Accuracy (\%) \\
\midrule
200 & 0.2 & 97.36 \\
200 & 0.2 & 97.14 \\
100 & 0.2 & 96.92 \\
200 & 0.1 & 96.70 \\
100 & 0.2 & 96.48 \\
50 & 0.2 & 96.26 \\
\bottomrule
\end{tabular}
\end{table}

Best Gradient Boosting setting: \texttt{n\_estimators=200}, \texttt{learning\_rate=0.2},
\texttt{max\_depth=3}. 5-fold CV F1: [0.983, 0.991, 0.965, 0.966, 0.991], mean = 0.9792. AdaBoost
edged out Gradient Boosting slightly on CV metrics, so AdaBoost is what represents "Boosting" in the
final comparison and plots later on, though both are reported here since the experiment asked for
Boosting more generally.

\begin{table}[H]
\centering
\caption{Stacked Ensemble -- Hyperparameter Evaluation (top configurations by CV accuracy)}
\begin{tabular}{ccc}
\toprule
SVM C & Decision Tree max\_depth & Avg CV Accuracy (\%) \\
\midrule
10 & 5 & 96.92 \\
10 & None & 96.70 \\
1 & 5 & 96.70 \\
1 & None & 96.48 \\
10 & 3 & 96.04 \\
1 & 3 & 95.82 \\
\bottomrule
\end{tabular}
\end{table}

Best Stacking setting: base SVM with \texttt{C=10}, Decision Tree with \texttt{max\_depth=5}, Naive
Bayes left at defaults, meta-learner Logistic Regression. 5-fold CV F1: [0.966, 0.982, 0.956, 0.983,
0.991], mean = 0.9755.

\section{Performance Comparison}
\begin{table}[H]
\centering
\caption{Performance Comparison of Ensemble Models (held-out test set)}
\begin{tabular}{lcccc}
\toprule
Model & Accuracy (\%) & Precision & Recall & F1 Score \\
\midrule
Bagging & 95.61 & 0.9589 & 0.9722 & 0.9655 \\
AdaBoost & 95.61 & 0.9467 & 0.9861 & 0.9660 \\
Gradient Boosting & 95.61 & 0.9467 & 0.9861 & 0.9660 \\
Stacked Ensemble & \textbf{97.37} & \textbf{0.9726} & 0.9861 & \textbf{0.9793} \\
\bottomrule
\end{tabular}
\end{table}

\section{Result Visualization}

\begin{figure}[H]
\centering
\includegraphics[width=\linewidth]{outputs/confusion_matrices_ensembles.png}
\caption{Confusion matrices on the held-out test set for Bagging, Boosting (AdaBoost), and the
Stacked Ensemble.}
\end{figure}

\begin{figure}[H]
\centering
\includegraphics[width=0.7\linewidth]{outputs/roc_curves_ensembles.png}
\caption{ROC curves comparing Bagging, AdaBoost and the Stacked Ensemble on the test set.}
\end{figure}

\textbf{Inference:} All three ensembles sit in a tight band above 95\% accuracy, so none of them is
struggling on this dataset, but Stacking is visibly the cleanest in the confusion matrix, with the
fewest misclassified malignant cases, which matters a lot here since a missed malignant case is the
expensive kind of mistake in a diagnostic setting. The ROC curves for all three hug the top-left
corner tightly, all comfortably above AUC = 0.98, so the gap between them is really in the harder
borderline cases rather than in broad separability.

\section{Observations}
\begin{itemize}
\item \textbf{How does Bagging reduce variance:} By training each Decision Tree on a different
bootstrap sample (here, half the rows and half the features per tree) and then averaging their votes,
the noisy, overfit quirks of any one tree get smoothed out. A single unpruned Decision Tree tends to
memorize the training data, but 50 of them voting together cancel out each other's individual mistakes.
\item \textbf{How does Boosting address bias:} Both AdaBoost and Gradient Boosting build models
sequentially, and each new one is trained specifically to fix what the previous ones got wrong, either
by reweighting misclassified samples (AdaBoost) or by fitting the residual errors (Gradient Boosting).
That's why Boosting tends to push a weak, high-bias base learner toward a much stronger fit, which
shows up here in AdaBoost's CV F1 of 0.984, the best of any single-family model tested.
\item \textbf{Why stacking benefits from heterogeneous models:} SVM, Naive Bayes and a Decision Tree
make mistakes in different, mostly uncorrelated ways -- SVM struggles where classes overlap in feature
space, Naive Bayes struggles when features aren't actually independent, and a Decision Tree struggles
with smooth boundaries. Feeding all three sets of predictions into a Logistic Regression meta-learner
lets it learn which base learner to trust in which situation, rather than betting everything on one
model's blind spots.
\item \textbf{Which ensemble performed best and why:} Stacking, at 97.37\% accuracy and 0.9793 F1,
came out ahead of both Bagging and Boosting on the test set. It's basically getting to have it both
ways -- SVM and Decision Tree cover different kinds of decision boundaries, Naive Bayes adds a cheap,
different-shaped opinion, and the meta-learner combines all three more intelligently than a simple
vote or weighted sum would.
\end{itemize}

\section{Discussion}
Bagging, AdaBoost and Gradient Boosting all landed at exactly the same 95.61\% test accuracy, which is
a bit of a coincidence but it does show that on a dataset this clean and well-behaved, different
ensemble philosophies can converge to a similar ceiling. AdaBoost and Gradient Boosting had identical
precision and recall too, though their actual cross-validation numbers (0.984 vs 0.979 mean F1) show
AdaBoost had a slight edge in training. Stacking's advantage on the test set lines up with its
somewhat lower raw CV F1 (0.9755) -- it wasn't necessarily the most "confident" model during
cross-validation, but its combination of decision boundaries generalized a little better to the actual
held-out data.

\section{Conclusion}
All three ensemble strategies clearly beat what a lone Decision Tree or a lone SVM would manage on
this dataset, but they get there differently: Bagging tames variance by averaging many decorrelated
trees, Boosting chases down bias by focusing each new learner on the previous ones' mistakes, and
Stacking blends genuinely different model types through a learned meta-combination. On the Breast
Cancer Wisconsin dataset specifically, Stacking came out on top, mainly because malignant and benign
cases here aren't separated by one single kind of boundary, so a mix of SVM, Naive Bayes and Decision
Tree opinions, properly weighted by the meta-learner, catches cases that any one of them alone would
miss.

\section{Learning Outcomes}
\begin{itemize}
\item Built and tuned a Bagging classifier, and saw directly how sampling fraction and estimator count
trade off against each other.
\item Compared two different Boosting algorithms (AdaBoost and Gradient Boosting) side by side rather
than treating "Boosting" as one single technique.
\item Practiced constructing a Stacked ensemble with genuinely different base learners and a proper
cross-validated meta-learner, instead of just averaging predictions.
\item Got a much more concrete sense of the bias-variance story behind ensembles: Bagging as a
variance-reduction tool, Boosting as a bias-reduction tool, Stacking as a way to combine different
error patterns.
\item Practiced reading ROC curves and confusion matrices specifically through a diagnostic lens,
where false negatives (missed malignant cases) matter more than the raw accuracy number suggests.
\end{itemize}

\section{References}
\begin{enumerate}
\item Scikit-learn documentation: \url{https://scikit-learn.org/stable/modules/ensemble.html}
\item UCI Machine Learning Repository, Breast Cancer Wisconsin (Diagnostic) Data Set:
\url{https://archive.ics.uci.edu/dataset/17/breast+cancer+wisconsin+diagnostic}
\item Ethem Alpaydin, ``Introduction to Machine Learning'', 3rd Edition, The MIT Press, 2014.
\end{enumerate}

\end{document}
